\chapter{Non-Abelian Gauge Theory as a Quantum Theory}

\begin{remarkbox}
QED is based on an Abelian gauge group. Non-Abelian gauge theory is richer: the
gauge fields themselves carry charge and interact with one another. This chapter
introduces the quantum structure behind Yang--Mills theory, preparing the way for
BRST symmetry and QCD.
\end{remarkbox}

\section{From Abelian to Non-Abelian Gauge Symmetry}

In QED, the gauge group is \(U(1)\). The gauge field is neutral, and the field
strength is linear in \(A_\mu\):
\[
    F_{\mu\nu}=\partial_\mu A_\nu-\partial_\nu A_\mu.
\]
In a non-Abelian gauge theory, the gauge group has noncommuting generators
\(T^a\):
\[
    [T^a,T^b]=if^{abc}T^c.
\]
The constants \(f^{abc}\) are the structure constants of the Lie algebra.

The gauge field carries an adjoint index:
\[
    A_\mu=A_\mu^aT^a.
\]
The covariant derivative is
\[
    D_\mu=\partial_\mu+igA_\mu^aT^a.
\]
Because the generators do not commute, the gauge field becomes self-interacting.

\section{Field Strength and Yang--Mills Action}

The non-Abelian field strength is defined by the commutator of covariant
derivatives:
\[
    [D_\mu,D_\nu]=igF_{\mu\nu}^aT^a.
\]
In components,
\[
    F_{\mu\nu}^a
    =\partial_\mu A_\nu^a-\partial_\nu A_\mu^a
      -g f^{abc}A_\mu^bA_\nu^c,
\]
up to sign conventions for the covariant derivative.

The Yang--Mills Lagrangian is
\[
    \mathcal{L}_{\mathrm{YM}}
    = -\frac{1}{4}F_{\mu\nu}^aF^{a\mu\nu}.
\]
Unlike the Maxwell Lagrangian, this contains interaction terms even without
matter. Expanding \(F_{\mu\nu}^aF^{a\mu\nu}\) gives cubic and quartic powers of
\(A_\mu^a\). Therefore gauge bosons interact with each other.

\section{Gauge Transformations}

A finite gauge transformation is
\[
    U(x)=e^{i\alpha^a(x)T^a}.
\]
Matter fields transform as
\[
    \psi(x)\longrightarrow U(x)\psi(x).
\]
The gauge field transforms as
\[
    A_\mu\longrightarrow U A_\mu U^{-1}+\frac{i}{g}U\partial_\mu U^{-1},
\]
with the precise sign depending on the convention for \(D_\mu\).

Infinitesimally, the gauge field changes by
\[
    \delta A_\mu^a = -D_\mu^{ab}\alpha^b,
\]
where the adjoint covariant derivative is
\[
    D_\mu^{ab}=\delta^{ab}\partial_\mu+g f^{acb}A_\mu^c.
\]
The nonlinearity of this transformation is another sign that non-Abelian gauge
theory is structurally different from QED.

\section{Gauge Fixing}

The path integral over \(A_\mu^a\) overcounts gauge-equivalent configurations.
As in the Abelian case, one must fix a gauge. A common covariant gauge-fixing
term is
\[
    \mathcal{L}_{\mathrm{gf}}
    = -\frac{1}{2\xi}(\partial_\mu A^{a\mu})^2.
\]
The gauge-fixed quadratic part of the action gives the gauge-boson propagator
\[
    D_{\mu\nu}^{ab}(k)
    = \frac{-i\delta^{ab}}{k^2+i\epsilon}
      \left(\eta_{\mu\nu}-(1-\xi)\frac{k_\mu k_\nu}{k^2+i\epsilon}\right).
\]
In Feynman gauge,
\[
    D_{\mu\nu}^{ab}(k)=\frac{-i\delta^{ab}\eta_{\mu\nu}}{k^2+i\epsilon}.
\]
The color index \(a\) is conserved along a free gauge-boson propagator.

\section{Faddeev--Popov Ghosts}

For non-Abelian gauge theory, the Faddeev--Popov determinant depends on the
gauge field. It is represented using ghost fields \(c^a,\bar c^a\). The ghost
Lagrangian in covariant gauge is
\[
    \mathcal{L}_{\mathrm{gh}}
    = \bar c^a\partial^\mu D_\mu^{ab}c^b.
\]
Expanding the covariant derivative gives a free ghost kinetic term and a
ghost--gauge-boson interaction.

Ghosts are anticommuting scalar fields. They are not physical particles and do
not appear as external states. They are required in covariant gauges to preserve
unitarity and gauge independence of physical amplitudes.

In Abelian QED, ghosts decouple in linear gauges. In non-Abelian theories, they
do not decouple because the determinant depends on \(A_\mu^a\).

\section{Gauge-Boson Self-Interactions}

The Yang--Mills field strength contains the nonlinear term
\[
    -g f^{abc}A_\mu^bA_\nu^c.
\]
Therefore the action contains cubic and quartic gauge-field interactions.
Diagrammatically, there are two new types of vertices:
\begin{itemize}
    \item a three-gauge-boson vertex proportional to \(g f^{abc}\);
    \item a four-gauge-boson vertex proportional to \(g^2\) and products of
    structure constants.
\end{itemize}
These vertices have no analogue in QED. They are responsible for many of the
special features of non-Abelian gauge theories, including asymptotic freedom.

The physical meaning is simple but profound: gauge bosons themselves carry the
charge associated with the gauge group.

\section{Matter Couplings}

Matter fields couple through the covariant derivative. For a Dirac fermion in a
representation \(R\), the Lagrangian is
\[
    \mathcal{L}_{\mathrm{matter}}
    = \bar\psi(i\gamma^\mu D_\mu-m)\psi.
\]
Expanding the derivative gives the interaction
\[
    \mathcal{L}_{\mathrm{int}}
    = -g\bar\psi\gamma^\mu T_R^a\psi A_\mu^a,
\]
up to sign conventions. The corresponding vertex factor is
\[
    -ig\gamma^\mu T_R^a.
\]
This is the non-Abelian analogue of the QED vertex \(-ie\gamma^\mu\), with the
charge \(e\) replaced by a matrix \(gT^a\).

The group generators track how matter transforms under the gauge group. In QCD,
quarks transform in the fundamental representation of \(SU(3)\), and gluons
transform in the adjoint representation.

\section{Feynman Rules: The Basic Set}

In a covariant gauge, the basic perturbative ingredients are:
\begin{itemize}
    \item gauge-boson propagator \(D_{\mu\nu}^{ab}(k)\);
    \item ghost propagator;
    \item fermion propagator, if matter is present;
    \item fermion--gauge-boson vertex;
    \item ghost--gauge-boson vertex;
    \item three-gauge-boson vertex;
    \item four-gauge-boson vertex.
\end{itemize}
Every diagram carries both Lorentz structure and group-theory structure.
Momentum conservation holds at each vertex, while group indices are contracted
according to the generators and structure constants.

The algebra is more complex than QED, but the logic is the same: propagators come
from quadratic terms, and vertices come from interaction terms in the
gauge-fixed action.

\section{Ward Identities and Slavnov--Taylor Identities}

In QED, gauge invariance leads to Ward identities. In non-Abelian gauge theory,
the corresponding identities are more complicated because gauge transformations
are nonlinear and ghosts are present. They are called Slavnov--Taylor identities.

These identities relate Green functions and ensure that unphysical polarizations
and ghosts cancel from physical observables. They also constrain
renormalization, ensuring that the gauge-fixed quantum theory preserves the
content of the original gauge symmetry.

BRST symmetry provides the most efficient modern derivation of these identities.
The next chapter develops BRST symmetry more systematically.

\section{Unitarity and Physical States}

Covariant quantization introduces unphysical gauge-field polarizations and ghost
fields. The full gauge-fixed state space is larger than the physical Hilbert
space. Physical states must be selected so that negative-norm and ghost degrees
of freedom do not appear as observable particles.

The cancellation of unphysical contributions is not diagram-by-diagram obvious.
It follows from the gauge symmetry encoded after gauge fixing, especially through
BRST symmetry. This is why ghosts, although unphysical, are essential in loop
calculations.

The physical \(S\)-matrix is unitary when restricted to the physical subspace.
This statement is one of the central consistency requirements of quantum gauge
theory.

\section{One-Loop Beta Function}

Non-Abelian gauge theories have a remarkable one-loop beta function. For an
\(SU(N)\) gauge theory with \(n_f\) Dirac fermions in the fundamental
representation, the leading beta function has the form
\[
    \beta(g)
    = -\frac{g^3}{16\pi^2}
      \left(\frac{11}{3}N-\frac{2}{3}n_f\right)+\cdots.
\]
The gauge-boson self-interactions give the positive \(11N/3\) contribution
inside the parentheses, leading to a negative beta function when \(n_f\) is not
too large.

This is the origin of asymptotic freedom. At high energies, the coupling becomes
small. At low energies, it grows, leading toward strong-coupling physics.

\section{Comparison with QED}

The contrast with QED is sharp. In QED, photons do not carry electric charge and
there are no photon self-interactions at the fundamental level. The leading beta
function is positive, and the coupling grows at high energies.

In non-Abelian gauge theory, gauge bosons carry the gauge charge and interact
with one another. Their contribution to the beta function can dominate over
matter screening and make the theory asymptotically free.

This difference is not a technical detail. It explains why QCD behaves very
differently from QED: quarks and gluons are weakly coupled at short distances but
strongly coupled at long distances.

\section{Toward QCD}

Quantum chromodynamics is the non-Abelian gauge theory with gauge group
\(SU(3)\). Its gauge bosons are gluons, and its matter fields are quarks. The
gluons carry color charge because they transform in the adjoint representation.

The Yang--Mills structure introduced in this chapter provides the foundation of
QCD:
\[
    \mathcal{L}_{\mathrm{QCD}}
    = -\frac{1}{4}G_{\mu\nu}^aG^{a\mu\nu}
      +\sum_f\bar q_f(i\gamma^\mu D_\mu-m_f)q_f.
\]
The later QCD chapter will build on this structure to discuss color, asymptotic
freedom, confinement, jets, and hadrons.

\begin{summarybox}
Non-Abelian gauge theories are quantum theories of gauge fields whose symmetry
generators do not commute. Their gauge bosons self-interact, ghosts do not
decouple in covariant gauges, and physical consistency is encoded in
Slavnov--Taylor identities and BRST symmetry. The one-loop beta function can be
negative, producing asymptotic freedom. This structure is the foundation of QCD
and of the gauge sector of the Standard Model.
\end{summarybox}

\section{Chapter Checklist}

After reading this chapter, one should be able to:
\begin{itemize}
    \item distinguish Abelian and non-Abelian gauge symmetry;
    \item write the non-Abelian field strength;
    \item explain why Yang--Mills theory contains gauge-boson self-interactions;
    \item write the covariant gauge-fixing term;
    \item explain why ghosts do not decouple in non-Abelian theories;
    \item identify the basic non-Abelian Feynman rules;
    \item write the matter--gauge-boson vertex;
    \item explain the role of Slavnov--Taylor identities;
    \item explain why BRST symmetry is needed for physical-state selection;
    \item state the qualitative origin of asymptotic freedom;
    \item connect non-Abelian gauge theory to QCD.
\end{itemize}
